\documentclass[10pt,a4paper]{amsart}
\usepackage{amsmath,amssymb,amsthm,graphicx}
\usepackage{verbatim}
\newenvironment{reference}{\comment}{\endcomment}
%\newenvironment{reference}{}{}
\newenvironment{slogan}{\comment}{\endcomment}
\newenvironment{history}{\comment}{\endcomment}

% For commutative diagrams we use Xy-pic
\usepackage[all]{xy}

% We use 2cell for 2-commutative diagrams.
\xyoption{2cell}
% Theorem environments.
%
\theoremstyle{plain}
\newtheorem{theorem}[subsection]{Theorem}
\newtheorem{proposition}[subsection]{Proposition}
\newtheorem{lemma}[subsection]{Lemma}

\theoremstyle{definition}
\newtheorem{definition}[subsection]{Definition}
\newtheorem{example}[subsection]{Example}
\newtheorem{exercise}[subsection]{Exercise}
\newtheorem{situation}[subsection]{Situation}

\theoremstyle{remark}
\newtheorem{remark}[subsection]{Remark}
\newtheorem{remarks}[subsection]{Remarks}

\numberwithin{equation}{subsection}

% Macros
%
\def\lim{\mathop{\mathrm{lim}}\nolimits}
\def\colim{\mathop{\mathrm{colim}}\nolimits}
\def\Spec{\mathop{\mathrm{Spec}}}
\def\Hom{\mathop{\mathrm{Hom}}\nolimits}
\def\Ext{\mathop{\mathrm{Ext}}\nolimits}
\def\SheafHom{\mathop{\mathcal{H}\!\mathit{om}}\nolimits}
\def\SheafExt{\mathop{\mathcal{E}\!\mathit{xt}}\nolimits}
\def\Sch{\mathit{Sch}}
\def\Mor{\mathop{\mathrm{Mor}}\nolimits}
\def\Ob{\mathop{\mathrm{Ob}}\nolimits}
\def\Sh{\mathop{\mathit{Sh}}\nolimits}
\def\NL{\mathop{N\!L}\nolimits}
\def\CH{\mathop{\mathrm{CH}}\nolimits}
\def\proetale{{pro\text{-}\acute{e}tale}}
\def\etale{{\acute{e}tale}}
\def\QCoh{\mathit{QCoh}}
\def\Ker{\mathop{\mathrm{Ker}}}
\def\Im{\mathop{\mathrm{Im}}}
\def\Coker{\mathop{\mathrm{Coker}}}
\def\Coim{\mathop{\mathrm{Coim}}}

% Boxtimes
%
\DeclareMathSymbol{\boxtimes}{\mathbin}{AMSa}{"02}

%
% Macros for moduli stacks/spaces
%
\def\QCohstack{\mathcal{QC}\!\mathit{oh}}
\def\Cohstack{\mathcal{C}\!\mathit{oh}}
\def\Spacesstack{\mathcal{S}\!\mathit{paces}}
\def\Quotfunctor{\mathrm{Quot}}
\def\Hilbfunctor{\mathrm{Hilb}}
\def\Curvesstack{\mathcal{C}\!\mathit{urves}}
\def\Polarizedstack{\mathcal{P}\!\mathit{olarized}}
\def\Complexesstack{\mathcal{C}\!\mathit{omplexes}}
% \Pic is the operator that assigns to X its picard group, usage \Pic(X)
% \Picardstack_{X/B} denotes the Picard stack of X over B
% \Picardfunctor_{X/B} denotes the Picard functor of X over B
\def\Pic{\mathop{\mathrm{Pic}}\nolimits}
\def\Picardstack{\mathcal{P}\!\mathit{ic}}
\def\Picardfunctor{\mathrm{Pic}}
\def\Deformationcategory{\mathcal{D}\!\mathit{ef}}
\usepackage{cleveref}
\crefname{originallemma}{original lemma}{original lemmas}
\Crefname{originallemma}{Original lemma}{Original lemmas}
\makeatletter
\newlabel{divisors-lemma-normal-effective-Cartier-divisor-S1}{{0B3S}{source}{}{sourcepointer-1}{}}
\newlabel{divisors-lemma-normal-effective-Cartier-divisor-S1@cref}{{[originallemma][0][]0B3S}{[1][0][]source}{}{}{}}
\makeatother
\crefname{originalreference}{original reference}{original references}
\Crefname{originalreference}{Original reference}{Original references}
\makeatletter
\newlabel{lemma-blow-up-normal-rational}{{0B4Z}{source}{}{sourcepointer-2}{}}
\newlabel{lemma-blow-up-normal-rational@cref}{{[originalreference][0][]0B4Z}{[1][0][]source}{}{}{}}
\makeatother
\crefname{originalreference}{original reference}{original references}
\Crefname{originalreference}{Original reference}{Original references}
\makeatletter
\newlabel{lemma-sections-powers-I-rational}{{0B4Y}{source}{}{sourcepointer-3}{}}
\newlabel{lemma-sections-powers-I-rational@cref}{{[originalreference][0][]0B4Y}{[1][0][]source}{}{}{}}
\makeatother
\crefname{originalreference}{original reference}{original references}
\Crefname{originalreference}{Original reference}{Original references}
\makeatletter
\newlabel{situation-rational}{{0B4W}{source}{}{sourcepointer-4}{}}
\newlabel{situation-rational@cref}{{[originalreference][0][]0B4W}{[1][0][]source}{}{}{}}
\makeatother
\crefname{originaldefinition}{original definition}{original definitions}
\Crefname{originaldefinition}{Original definition}{Original definitions}
\makeatletter
\newlabel{varieties-definition-degree-invertible-sheaf}{{0AYR}{source}{}{sourcepointer-5}{}}
\newlabel{varieties-definition-degree-invertible-sheaf@cref}{{[originaldefinition][0][]0AYR}{[1][0][]source}{}{}{}}
\makeatother
\title{Cohomology of the exceptional curve of the blowup in the rational situation}
\author{The Stacks Project Authors}
\date{stacks-project:resolve.tex:local-snapshot-20261012; GNU FDL 1.2 or later}
\begin{document}
\maketitle
\noindent\textit{Editorial note.} All delivered mathematical text is the author's own, copied from the pinned Stacks Project chapter file resolve.tex; no statement or proof step was rewritten, reordered or removed.\par\medskip
\noindent\textit{Editorial note.} The counted proof is the author's complete argument: the very ampleness and Cohen-Macaulayness statements, the cohomological computation through the comparison of powers, and the degree computation giving the regularity criterion.\par\medskip
\noindent\textit{Editorial note.} The label lemma-cohomology-blow-up-rational names the counted statement, and the proof environment immediately following it in the source is delivered verbatim as its proof.\par\medskip
\noindent\textit{Editorial note.} The counted statement is carried by the author's own lemma environment in the chapter source.\par\medskip
\noindent\textit{Editorial note.} Pre-open gate: python operations/agent-authoring/tools/precard\_gate.py --source resolve.tex --label lemma-cohomology-blow-up-rational --start 2580 --end 2615 --local-id 28110 returned RESULT=OK (L28110 free, occupancy FREE, public interval clear). Pre-assembly tag self-check: canonical tag resolved and not present in reference/public-items.jsonl.\par\medskip
\noindent\textit{Editorial note.} The statement and proof are copied from the pinned chapter file; the cited inputs are kept as external pointers with their canonical Stacks tags.\par\medskip
\noindent\textit{Allowed standard prerequisites (editorial).}\begin{itemize}
\item The blowup is normal in the rational situation. Same chapter, the lemma that the blowup in the rational situation is normal. Makes the exceptional divisor an effective Cartier divisor on a normal scheme, hence Cohen-Macaulay.
\item Sections and vanishing for the powers of the ideal sheaf of the exceptional divisor. Same chapter, the computation of the sections of the powers of the ideal sheaf and the vanishing of their first cohomology. Supplies the two cohomological statements through the short exact sequences comparing consecutive powers.
\item Cohen-Macaulayness of an effective Cartier divisor on a normal scheme. Divisors, the lemma that a normal effective Cartier divisor satisfies the first Serre condition. Gives that the exceptional curve is Cohen-Macaulay.
\item Degree of an invertible sheaf on a curve through the Euler characteristic. Varieties, the definition of the degree of an invertible sheaf. Turns the difference of Euler characteristics into the degree and hence into the embedding dimension of the maximal ideal minus one.
\end{itemize}\medskip
\ifdefined\hypertarget\hypertarget{sourcepointer-1}{}\fi
\noindent\textit{Source reference.} Original lemma 0B3S: sources/stacks-chow/tags.txt, source lines 13975--13975. Prerequisite pointer invoked at proof line(s) 2602; it names a standard fact of a sibling chapter that the counted argument uses, and every step of the counted proof built on it is retained verbatim in the delivered body.\par\medskip
\ifdefined\hypertarget\hypertarget{sourcepointer-2}{}\fi
\noindent\textit{Source reference.} Original reference 0B4Z: sources/stacks-resolve/pinned-original-source.tex, source lines 2534--2534. Prerequisite pointer invoked at proof line(s) 2600; it names a standard fact of the same chapter that the counted argument uses, and every step of the counted proof built on it is retained verbatim in the delivered body.\par\medskip
\ifdefined\hypertarget\hypertarget{sourcepointer-3}{}\fi
\noindent\textit{Source reference.} Original reference 0B4Y: sources/stacks-resolve/pinned-original-source.tex, source lines 2488--2488. Prerequisite pointer invoked at proof line(s) 2603; it names a standard fact of the same chapter that the counted argument uses, and every step of the counted proof built on it is retained verbatim in the delivered body.\par\medskip
\ifdefined\hypertarget\hypertarget{sourcepointer-4}{}\fi
\noindent\textit{Source reference.} Original reference 0B4W: sources/stacks-resolve/pinned-original-source.tex, source lines 2455--2455. Prerequisite pointer invoked at proof line(s) 2582; it names a standard fact of the same chapter that the counted argument uses, and every step of the counted proof built on it is retained verbatim in the delivered body.\par\medskip
\ifdefined\hypertarget\hypertarget{sourcepointer-5}{}\fi
\noindent\textit{Source reference.} Original definition 0AYR: sources/stacks-chow/tags.txt, source lines 13799--13799. Prerequisite pointer invoked at proof line(s) 2613; it names a standard fact of a sibling chapter that the counted argument uses, and every step of the counted proof built on it is retained verbatim in the delivered body.\par\medskip
\begin{lemma}
\label{lemma-cohomology-blow-up-rational}
In Situation \ref{situation-rational}.
Let $X$ be the blowup of $\Spec(A)$ in $\mathfrak m$. Let $E \subset X$
be the exceptional divisor. With $\mathcal{O}_X(1) = \mathcal{I}$ as
usual and $\mathcal{O}_E(1) = \mathcal{O}_X(1)|_E$ we have
\begin{enumerate}
\item $E$ is a proper Cohen-Macaulay curve over $\kappa$.
\item $\mathcal{O}_E(1)$ is very ample
\item $\deg(\mathcal{O}_E(1)) \geq 1$ and equality holds only if
$A$ is a regular local ring,
\item $H^1(E, \mathcal{O}_E(n)) = 0$ for $n \geq 0$, and
\item $H^0(E, \mathcal{O}_E(n)) = \mathfrak m^n/\mathfrak m^{n + 1}$
for $n \geq 0$.
\end{enumerate}
\end{lemma}

\begin{proof}
Since $\mathcal{O}_X(1)$ is very ample by construction, we see that
its restriction to the special fibre $E$ is very ample as well.
By Lemma \ref{lemma-blow-up-normal-rational} the scheme $X$ is normal.
Then $E$ is Cohen-Macaulay by
Divisors, Lemma \ref{divisors-lemma-normal-effective-Cartier-divisor-S1}.
Lemma \ref{lemma-sections-powers-I-rational} applies and we obtain
(4) and (5) from the exact sequences
$$
0 \to \mathcal{I}^{n + 1} \to \mathcal{I}^n \to i_*\mathcal{O}_E(n) \to 0
$$
and the long exact cohomology sequence. In particular, we see that
$$
\deg(\mathcal{O}_E(1)) = \chi(E, \mathcal{O}_E(1)) - \chi(E, \mathcal{O}_E) =
\dim(\mathfrak m/\mathfrak m^2) - 1
$$
by Varieties, Definition \ref{varieties-definition-degree-invertible-sheaf}.
Thus (3) follows as well.
\end{proof}

\end{document}
